\documentclass[11pt]{article}

\usepackage[a4paper,margin=1in]{geometry}
\usepackage{amsmath,amssymb,amsthm,mathtools,bm}
\usepackage{booktabs}
\usepackage{enumitem}
\usepackage{algorithm}
\usepackage{algpseudocode}
\usepackage{microtype}
\usepackage[dvipsnames]{xcolor}
\usepackage[colorlinks=true,linkcolor=NavyBlue,citecolor=NavyBlue,urlcolor=NavyBlue]{hyperref}

% ---- math shortcuts ----
\DeclareMathOperator*{\argmax}{arg\,max}
\DeclareMathOperator*{\argmin}{arg\,min}
\newcommand{\E}{\mathbb{E}}
\newcommand{\R}{\mathbb{R}}
\newcommand{\Teams}{\mathcal{T}}
\newcommand{\Probs}{\mathcal{P}}
\newcommand{\Conts}{\mathcal{C}}
\newcommand{\Obs}{\mathcal{O}}
\newcommand{\given}{\,\vert\,}
\newcommand{\indic}[1]{\mathbf{1}\!\left[#1\right]}

\theoremstyle{definition}
\newtheorem{definition}{Definition}
\newtheorem{remark}{Remark}

\title{\textbf{Rating ICPC Problems from Contest Standings}\\[2pt]
\large A two-layer model with cross-region standardization}
\author{Bilguun}
\date{\today}

\begin{document}
\maketitle

\begin{abstract}
\noindent
Competitive programming judges such as Codeforces and AtCoder assign every
contestant a rating, which lets tools like \texttt{clist.by} estimate a
problem's difficulty by a single Elo inversion: find the rating level at which
the expected number of solvers matches the observed number. ICPC contests
provide no such native rating, so the contestant--strength layer must be built
from the standings themselves. This note develops two estimators for ICPC
problem difficulty from standings data (teams, solved problems, and per-problem
solve times) and shows how teams that recur across regional contests, the World
Finals, and the Universal Cup act as a linking signal that places problems from
different regions onto a single comparable scale.
\end{abstract}

\section{Setup and notation}\label{sec:setup}

Let $\Teams$ be the set of competitors (teams, or individuals; see
Remark~\ref{rem:identity}), $\Probs$ the set of problems, and $\Conts$ the set
of contests. Each contest $c \in \Conts$ has a duration $T_c$ and a fixed
problem set $\Probs_c \subseteq \Probs$. Write $\Teams_c$ for the teams that
participate in $c$. The observed data are
\begin{itemize}[leftmargin=1.4em,itemsep=2pt]
  \item $y_{tp} \in \{0,1\}$: whether team $t$ solved problem $p$
        (during the contest in which $p$ appeared);
  \item $\tau_{tp} \in [0,T_c]$: the time at which $t$ solved $p$, defined only
        when $y_{tp}=1$;
  \item $r_{t,c}$: the final rank of team $t$ in contest $c$.
\end{itemize}
We estimate two latent quantities:
\[
  \theta_t \in \R \quad\text{(ability of team $t$)},
  \qquad
  b_p \in \R \quad\text{(difficulty of problem $p$)} .
\]
Both live on the same scale, anchored as in Section~\ref{sec:link}.
Let $\Obs = \{(t,p) : t \text{ attempted the contest containing } p\}$ be the
set of competitor--problem pairs that are actually informative.

\section{The Elo inversion primitive}\label{sec:elo}

Throughout we reuse the Codeforces logistic convention, in which a competitor of
ability $\theta$ solves a problem of difficulty $b$ with probability
\begin{equation}\label{eq:pi}
  \pi(\theta,b) \;=\; \frac{1}{1 + 10^{\,(b-\theta)/400}}
  \;=\; \sigma\!\Big(\tfrac{\theta - b}{s}\Big),
  \qquad s = \frac{400}{\ln 10} \approx 173.7,
\end{equation}
where $\sigma(x) = (1+e^{-x})^{-1}$. The scale is calibrated so that
$\pi=\tfrac12$ when $\theta=b$, $\pi\approx0.76$ at a $200$-point edge, and
$\pi\approx0.91$ at a $400$-point edge. The pairwise win probability between two
competitors uses the same form, $E(\theta,\phi)=\pi(\theta,\phi)$.

The single primitive underlying everything below inverts a weighted sum of
\eqref{eq:pi} against a target. Given competitors with abilities
$\{\theta_t\}_{t\in I}$ and nonnegative weights $\{w_t\}$, define

\begin{definition}[Weighted rating]\label{def:wr}
\[
  \mathrm{WR}\big(\{(w_t,\theta_t)\}_{t\in I},\, S\big)
  \;=\;
  \text{the unique } b \text{ such that }
  \sum_{t \in I} w_t\, \pi(\theta_t, b) \;=\; S .
\]
\end{definition}

\noindent
The left-hand side is strictly decreasing and continuous in $b$, so the root is
unique and is found by bisection. Two specializations are used repeatedly:
\begin{align}
  \text{problem difficulty:}\quad
    & b_p = \mathrm{WR}\big(\{(w_t,\theta_t)\}_{t\in\Teams_c},\, S_p\big),
      \quad S_p = \sum_{t\in\Teams_c} w_t\, y_{tp};
      \label{eq:bp}\\[2pt]
  \text{performance rating:}\quad
    & \rho_{t,c} : \;\; 1 + \!\!\sum_{u\in\Teams_c,\,u\neq t}\!\! \pi(\theta_u,\rho_{t,c}) = r_{t,c}.
      \label{eq:perf}
\end{align}
Equation~\eqref{eq:bp} is exactly the difficulty estimate computed by
\texttt{clist.by}; the only difference in our setting is that
$\{\theta_t\}$ is unknown and must be estimated jointly
(Sections~\ref{sec:fixedpoint}--\ref{sec:irt}). Equation~\eqref{eq:perf} reads
``the ability whose expected rank equals the achieved rank.''

\paragraph{Experience weighting and team composition.}
Following the \texttt{clist} heuristic, a competitor seen in few prior contests
contributes less, so its unsettled ability cannot distort a difficulty estimate:
\begin{equation}\label{eq:weight}
  w_t = 1 - 0.9^{\,n_t + 1},
\end{equation}
where $n_t$ is the number of prior contests of $t$. When a team's strength must
be summarized from its members' abilities $\{\theta_m\}_{m\in t}$, use the same
primitive at the median quantile,
\begin{equation}\label{eq:team}
  \theta_t = \mathrm{WR}\big(\{(1,\theta_m)\}_{m\in t},\, \tfrac12\big).
\end{equation}

\section{Architecture A: alternating fixed point}\label{sec:fixedpoint}

The first estimator mirrors \texttt{clist} but adds the ability layer that
\texttt{clist} obtains for free from the host judge. Abilities and difficulties
are mutually dependent, so they are solved by alternation. The coupling across
contests happens entirely in the ability update~\eqref{eq:update}: a team that
played several contests gets one ability blending all of them, and that single
value then pulls every one of those contests' difficulty estimates toward a
common scale.

\begin{equation}\label{eq:update}
  \theta_t \;\leftarrow\;
  \frac{\displaystyle\sum_{c\,:\,t\in\Teams_c} w_{t,c}\,\tfrac{1}{2}\big(\rho_{t,c} + \theta_t^{\text{prior}}\big)}
       {\displaystyle\sum_{c\,:\,t\in\Teams_c} w_{t,c}} .
\end{equation}

\begin{algorithm}[h]
\caption{Alternating estimation of abilities and difficulties}
\label{alg:fp}
\begin{algorithmic}[1]
\State initialize $\theta_t \gets \theta_t^{\text{prior}}$ for all $t$ \Comment{from CF bootstrap, \eqref{eq:cfprior}}
\Repeat
  \For{each contest $c \in \Conts$}
    \For{each team $t \in \Teams_c$}
      \State $\rho_{t,c} \gets$ performance rating by \eqref{eq:perf}
    \EndFor
  \EndFor
  \For{each team $t \in \Teams$}
    \State update $\theta_t$ by \eqref{eq:update}
  \EndFor
\Until{$\max_t |\Delta\theta_t| < \varepsilon$}
\For{each problem $p$ in contest $c$}
  \State $b_p \gets \mathrm{WR}\big(\{(w_t,\theta_t)\}_{t\in\Teams_c},\, S_p\big)$ by \eqref{eq:bp}
\EndFor
\end{algorithmic}
\end{algorithm}

\noindent
Algorithm~\ref{alg:fp} is cheap, reuses the \texttt{clist} code almost verbatim,
and yields a usable first version. Its weaknesses are the ones \texttt{clist}
inherits: it returns point estimates with no uncertainty, and a problem with very
few in-contest solvers can be badly mis-rated when a single strong-but-low-rated
team solves it. These motivate Architecture~B.

\section{Architecture B: joint item-response model}\label{sec:irt}

Treat each solve as a binary response governed by the gap between ability and
difficulty. This is the Rasch model, the item-response cousin of
Bradley--Terry:
\begin{equation}\label{eq:rasch}
  \Pr(y_{tp}=1 \given \theta_t, b_p)
  \;=\; \sigma\!\Big(\tfrac{\theta_t - b_p}{s}\Big)
  \;=\; \pi(\theta_t, b_p).
\end{equation}
A two-parameter (2PL) extension adds a per-problem discrimination $a_p>0$,
allowing some problems to separate strong from weak teams more sharply than
others:
\begin{equation}\label{eq:twopl}
  \Pr(y_{tp}=1 \given \theta_t, b_p, a_p)
  \;=\; \sigma\!\big(a_p(\theta_t - b_p)\big).
\end{equation}
Assuming conditional independence of responses given the latent parameters, the
log-likelihood over all observed pairs is
\begin{equation}\label{eq:loglik}
  \ell(\bm\theta,\bm b)
  = \sum_{(t,p)\in\Obs}
    \Big[\, y_{tp}\log \pi(\theta_t,b_p)
      + (1-y_{tp})\log\big(1-\pi(\theta_t,b_p)\big) \Big].
\end{equation}
The same $\theta_t$ appears in the likelihood terms of every contest team $t$
entered, so contests sharing teams are linked automatically; no post-hoc
rescaling is needed.

\paragraph{Priors and the MAP objective.}
Gaussian priors regularize sparsely observed teams and problems---crucial for
small regional fields:
\begin{equation}\label{eq:priors}
  \theta_t \sim \mathcal{N}\big(\mu_t,\ \sigma_\theta^2\big),
  \qquad
  b_p \sim \mathcal{N}\big(\mu_b,\ \sigma_b^2\big),
\end{equation}
with the team prior mean $\mu_t$ supplied by the bootstrap of
Section~\ref{sec:link}. The estimate is the maximum-a-posteriori point (or the
full posterior, via MCMC / variational inference):
\begin{equation}\label{eq:map}
  (\hat{\bm\theta}, \hat{\bm b})
  = \argmax_{\bm\theta,\bm b}\;
    \ell(\bm\theta,\bm b)
    - \frac{1}{2\sigma_\theta^2}\sum_t (\theta_t-\mu_t)^2
    - \frac{1}{2\sigma_b^2}\sum_p (b_p-\mu_b)^2 .
\end{equation}
A problem solved by nobody contributes only $\sum_t \log(1-\pi(\theta_t,b_p))$,
which pushes $b_p$ upward but is held finite by the prior---so the pathological
$b_p\to+\infty$ of a bare likelihood does not occur.

\section{Using solve times: a survival formulation}\label{sec:survival}

The per-problem solve time $\tau_{tp}$ is signal that \texttt{clist} discards.
Model the act of solving as a constant-hazard (exponential) process over the
contest window, with proportional hazards in the ability--difficulty gap:
\begin{equation}\label{eq:hazard}
  \lambda_{tp} = \lambda_0 \exp\!\Big(\tfrac{\theta_t - b_p}{s}\Big).
\end{equation}
A team that solves $p$ at time $\tau_{tp}$ contributes the event density, and a
team that never solves $p$ contributes the survival probability to the contest
end (a right-censored observation). The per-pair log-likelihood is
\begin{equation}\label{eq:survlik}
  \ell_{tp} =
  \begin{cases}
    \log \lambda_{tp} - \lambda_{tp}\,\tau_{tp}, & y_{tp}=1 \ \text{(solved at } \tau_{tp}),\\[4pt]
    -\,\lambda_{tp}\, T_c, & y_{tp}=0 \ \text{(censored at } T_c).
  \end{cases}
\end{equation}
This single likelihood uses the solve-time distribution \emph{and} handles
zero-solve problems coherently, recovering the binary model of
Section~\ref{sec:irt} as the special case where only the order of
$\{y_{tp}\}$ is retained.

\begin{remark}
Contest length enters explicitly through $T_c$, which partly corrects the bias
from comparing a 5-hour regional against a 2-hour short-format round.
\end{remark}

\section{Cross-region standardization}\label{sec:link}

A single contest identifies abilities and difficulties only \emph{relative} to
each other: since \eqref{eq:rasch} depends on $\theta_t - b_p$, adding a constant
$\kappa$ to every $\theta_t$ and every $b_p$ leaves the likelihood unchanged.
Within one contest this is a single global shift. Across contests with disjoint
teams \emph{and} disjoint problems there is no shared parameter at all, so their
scales float independently and their difficulties are not comparable.

\paragraph{The linking graph.}
Form the graph $G=(\Conts, E)$ whose vertices are contests and whose edges are
\begin{equation}\label{eq:graph}
  (c,c') \in E \iff |\Teams_c \cap \Teams_{c'}| \ge k ,
\end{equation}
for a small threshold $k$. Two contests are on a common scale precisely when they
lie in the same connected component of $G$: a shared team carries one ability into
both contests' likelihoods, fixing their relative shift. The World Finals (one
problem set faced by teams qualifying from many regions) and the Universal Cup
(recurring shared rounds across regions) are the high-degree hubs that merge the
otherwise isolated regional components. In practice the Universal Cup is the
stronger linker, contributing far more cross-region team pairs than the
once-a-year Finals.

\paragraph{Anchoring.}
Within a connected component, one global degree of freedom remains. Fix it by an
identifiability constraint---either centering,
\begin{equation}\label{eq:center}
  \frac{1}{|\Teams|}\sum_{t\in\Teams}\theta_t = 0,
\end{equation}
or, preferably, by anchoring to a familiar external scale through the Codeforces
ratings of team members. For each team with members carrying Codeforces ratings
$\{R_m\}_{m\in t}$, set the prior mean by composition~\eqref{eq:team},
\begin{equation}\label{eq:cfprior}
  \mu_t = \mathrm{WR}\big(\{(1,R_m)\}_{m\in t},\, \tfrac12\big).
\end{equation}
This simultaneously solves cold-start (teams with no contest history get a
sensible prior), removes the global shift, and renders every output
interpretable in Codeforces-equivalent points.

\section{Practical notes}\label{sec:notes}

\begin{remark}[Competitor identity]\label{rem:identity}
A team's roster changes between seasons, so a team label is not a stable
identity across years. Where rosters are known, modeling at the individual level
and composing team strength by \eqref{eq:team} is preferable: it is more stable
and it densifies the linking graph, since a single competitor recurs across
seasons and appears in the Universal Cup with different teammates.
\end{remark}

\begin{remark}[Entity resolution]
The linking graph is only as good as the identifiers feeding it. If a team or
competitor is spelled inconsistently across regional standings, the Finals, the
Universal Cup, and mirror sources, the linking edge silently disappears and the
standardization quietly degrades. Identifier reconciliation is a prerequisite,
not an afterthought.
\end{remark}

\begin{remark}[Recommended path]
Ship Architecture~A first, bootstrapped by \eqref{eq:cfprior} and validated
against problems with an external difficulty opinion (Finals problems with known
editorial difficulty, or problems mirrored on a rated judge). Then move the final
standardized ratings to the Bayesian model of Section~\ref{sec:irt} with the
solve-time likelihood of Section~\ref{sec:survival}, once the linking graph's
connectivity is established.
\end{remark}

\section*{Pointers}
\addcontentsline{toc}{section}{Pointers}
\small
The Elo inversion and the $0.5/200/400$ calibration follow the Codeforces rating
and problem-difficulty conventions; the difficulty estimator~\eqref{eq:bp},
experience weight~\eqref{eq:weight}, and team composition~\eqref{eq:team} follow
the open-source \texttt{calculate\_problem\_rating} logic in
\href{https://github.com/aropan/clist}{\texttt{github.com/aropan/clist}}. The
item-response treatment is standard Rasch / 2PL; the survival formulation is a
proportional-hazards model with right-censoring.

\end{document}